\lecture{7.8}{Doubly Robust Inference with Flexible Learners}
\label{lec:7-8}

\subsection{Flexible Nuisance Estimation}

Flexible learners are attractive because they reduce reliance on simple parametric
models for $\Qbar$ and $g$. However, flexible nuisance estimation complicates
inference. Consistency of AIPTW or TMLE can hold if either nuisance function is
consistently estimated, but standard Wald inference usually requires more: a tractable
asymptotic linear expansion with consistently estimable variance.

When both nuisance functions are consistently estimated at adequate rates, standard
AIPTW and standard TMLE often have valid influence-function-based inference. The
harder question is what happens when only one nuisance function is consistently
estimated and flexible methods are used.

\subsection{Standard TMLE}

Standard TMLE targets the efficient influence function equation
\[
\frac{1}{n}\sum_{i=1}^nD_n^*(O_i)\approx 0,
\]
where $D_n^*$ is the efficient influence function with $\Qbar_n^*$ and $g_n$ plugged
in. This targeting step gives TMLE the same first-order expansion as AIPTW when the
nuisance estimators behave well.

Standard TMLE is doubly robust for consistency: if either $\Qbar_n$ or $g_n$ is
consistent, the parameter estimate can still be consistent. But this does not
automatically guarantee valid confidence intervals when the other nuisance estimator
is inconsistent and flexible estimation is used.

\subsection{TMLE for Doubly Robust Inference}

Doubly robust inference means more than doubly robust consistency. It means the
estimator has a usable asymptotic normal distribution and variance estimator even
when only one nuisance component is consistently estimated.

Specialized TMLEs can be constructed to target additional equations beyond the usual
efficient influence-function equation. These methods may update both the outcome
regression and the propensity score, often iteratively, to remove additional second-order
terms that would otherwise prevent valid inference.

The technical details are advanced, but the practical message is clear: standard AIPTW
and standard TMLE should not be assumed to deliver valid inference under arbitrary
flexible nuisance estimation if one nuisance function is badly inconsistent.

\subsection{Sample-Size Effects}

Finite-sample performance depends strongly on sample size. In small samples,
flexible learners can overfit, estimated propensity scores can be unstable, and
influence-function approximations can be poor. As sample size grows, flexible
regression can better approximate the nuisance functions, and coverage of
influence-function intervals often improves.

The rate at which coverage improves depends on the complexity of the true nuisance
functions, the learner library, positivity, and the amount of outcome noise.

\subsection{Coverage Probability}

Coverage probability is the repeated-sampling probability that a confidence interval
contains the true parameter. A nominal 95\% interval has good coverage if, across many
repeated samples, it contains the target about 95\% of the time.

Simulation studies comparing naive g-computation, standard TMLE, and TMLE designed
for doubly robust inference often show that coverage can be poor at small sample
sizes, especially when one nuisance regression is misspecified. Methods designed for
doubly robust inference can improve coverage in such settings, but they are more
complex and require careful implementation.

\subsection{Practical Cautions}

\begin{estimationbox}
Practical guidance for flexible doubly robust estimation:
\begin{itemize}
\item Use rich but pre-specified nuisance-function libraries.
\item Inspect estimated propensity scores and clever covariates for extreme values.
\item Use cross-fitting or sample splitting when required by the estimator's theory.
\item Report whether inference relies on both nuisance functions being consistently
estimated or on a method designed for doubly robust inference.
\item Treat small-sample Wald intervals cautiously when weights or influence-function
values are highly variable.
\end{itemize}
\end{estimationbox}

Flexible learners are valuable, but they do not make inference automatic. The more
adaptive the nuisance estimation procedure, the more important it is to use an
estimator and inference procedure whose theoretical conditions match the analysis.

\begin{keyidea}
Doubly robust consistency says that the point estimator can remain consistent if one
nuisance function is correct. Doubly robust inference is stronger: it asks for valid
standard errors and confidence intervals under that same one-correct-nuisance
scenario.
\end{keyidea}

\begin{assumptionbox}
No amount of flexible nuisance estimation can diagnose unmeasured confounding from
the observed data alone. Doubly robust inference concerns statistical robustness after
identification assumptions have been accepted.
\end{assumptionbox}
